\section{Scope, contributions, and non-claims}
\label{sec:scope}

This paper develops a \emph{deterministic substrate theory fragment} with an explicit goal:
to formulate a coordinate-free continuum model (using an isotropic constitutive baseline) with clear points of contact to simulations and diagnostics,
and to make the remaining open problems visible.
A \emph{cubic lattice} is treated as a candidate \emph{microphysical} substrate, while the continuum formulation is
retained as the \emph{long-wavelength effective description} and as a compact way to express isotropic constitutive
baselines.

\subsection{Positioning relative to established physics}
The Standard Model and general relativity are treated here as \emph{empirically validated effective frameworks}.
Nothing in this paper is intended to reduce confidence in their experimental domain of applicability.
The present work instead targets a different explanatory layer: it asks whether familiar structures such as
effective Lorentz symmetry, gauge connections, and quantized particle spectra could arise as emergent regularities
from a smaller set of ground ingredients.

\subsection{Core building blocks}
The theory fragment developed here is organized around six coupled elements:
\begin{enumerate}
\item \textbf{Lattice substrate and continuum limit.}
A cubic lattice of nodes/cells embedded in $\R^4$ is taken as a candidate microphysical substrate.
For wavelengths large compared to the lattice spacing, the dynamics admits a continuum approximation by an
isotropic hyperelastic action on an intrinsic $3$-domain $\Omega\subset\R^3$ (\Cref{sec:continuum-model}).

\item \textbf{Continuum 3D brane in 4D embedding (effective description).}
A local hyperelastic action on an intrinsic $3$-manifold $\Omega\subset\R^3$ (taken in an isotropic baseline form, i.e. without built-in preferred directions) determines the
dynamics of an embedding $\vecX:\Omega\times\R\to\R^4$ (\Cref{sec:continuum-model}).

\item \textbf{Gravity as contraction from transverse excitation.}
Because the brane is embedded, transverse gradients in $X^4$ change intrinsic distances through the induced metric.
Under homogeneous pre-tension ($\alpha<1$), localized transverse excitation therefore produces an inward lateral pull
(contraction) and a slowly varying background strain field. In a weak, slowly varying regime we interpret this
transverse--lateral backreaction as the qualitative origin of a long-range gravity-like channel
(\Cref{subsec:gravity-channel}).

\item \textbf{Ordered narrowband carrier sector.}
We assume the substrate admits an ordered, narrowband carrier sector (taken to be primarily longitudinal)
that provides phase rigidity and enables a controlled multiscale envelope description
(\Cref{sec:geophase,sec:effective-field}).

\item \textbf{Berry/Wilczek--Zee connections as gauge structure.}
Adiabatic transport of a polarization subspace induces a Berry connection (rank-1) or a Wilczek--Zee
connection (non-Abelian).
These connections are interpreted as the natural origin of gauge degrees of freedom in a coarse-grained
description, and provide the bridge to a $U(1)$-like gauge picture and its non-Abelian extension
(\Cref{sec:geophase,sec:effective-field}).

\item \textbf{Axis-aligned triplet and coupling-induced mixing (toward $SU(3)$).}
On a cubic substrate, an axis-aligned narrowband triplet $\Psi=(\psi_x,\psi_y,\psi_z)^\top$ is natural.
For $g_{\mathrm{mix}}=0$ the three phases behave as $U(1)\times U(1)\times U(1)$; for $g_{\mathrm{mix}}>0$ eigenmodes mix and the correct
adiabatic transport is a $U(3)$ Wilczek--Zee connection (with a traceless $SU(3)$ part).

\item \textbf{Solitonic particles and a Dirac-envelope split.}
Particle-like states are localized solitons whose spatial structure is organized by spherical symmetry and
spherical-harmonic mode families (\Cref{sec:localized}).
Conceptually, we separate (i) particle-internal structure (topology/geometry and its internal labels, including
matter/antimatter distinction and spin state) from (ii) the underlying local phase structure of the carrier sector.
The standard Dirac equation mixes both; here it is treated as an effective envelope description that couples the
particle sector to the local phase/gauge structure.
The spin-$\tfrac12$ $4\pi$ aspect is attributed to geometric phase/holonomy of an internal polarization subspace.
\end{enumerate}

\subsection{Non-claims (what is \emph{not} delivered here)}
To keep the scope honest, several items are explicitly \emph{not} claimed in this version:
\begin{itemize}
\item a derivation of the full Standard Model gauge group, generations, or precise coupling constants,
\item a completed quantitative reduction of the contraction channel to a Newtonian limit with a derived value of $G$,
\item a complete derivation of Maxwell/Yang--Mills dynamics for the emergent connection in all regimes,
\item a closed multi-particle theory including scattering cross sections, bound-state spectra, and entanglement
phenomenology.
\end{itemize}

\subsection{Falsifiable signatures and diagnostic targets}
The framework makes concrete predictions that can be tested (and potentially ruled out) in numerical experiments on
the lattice substrate and in any future phenomenological matching.
Key diagnostic targets include:
\begin{itemize}
\item \textbf{Lorentz-violation proxies:} direction-dependent wave speeds, birefringence, and dispersion onset as functions of wavelength relative to lattice spacing.
\item \textbf{Branch leakage / adiabaticity failure:} measurable conversion between polarization branches when the gap condition \Cref{eq:adiabatic-condition} is violated.
\item \textbf{Gauge-holonomy invariants:} gauge-invariant loop observables (rank-1 phase; rank-$n$ holonomy spectrum) extracted by the protocol in \Cref{subsec:diagnostic-protocol}.
\item \textbf{Non-Abelianity:} noncommuting holonomies for distinct loops in regimes where a multi-dimensional transported subspace is retained.
\item \textbf{Gravity-channel response:} a slow contraction/dilation field sourced by transverse excitation and its influence on test-wave propagation in weak, slowly varying limits.
\item \textbf{Localized-mode stability:} long-time persistence and perturbation response of soliton-like states (or clear failure modes if they do not exist in the chosen constitutive/stencil class).
\end{itemize}
The contribution is therefore a well-specified substrate model plus concrete diagnostics and simulation checks that can
falsify or refine the proposed emergence mechanisms.
